% ch07_u08 —— 原书 p.306–p.315（7.11.7–7.11.15(ii)）

\SourcePage{306}

以下不再提 D-晶体。本书正文中我们使用通常的 D-模（我们遇到的归纳概形是仿射 Grassmann 概形，它们形式光滑）。不过要注意，正是 D-晶体的进路让下述事实成为显然（当 $Y$ 是 Schubert 胞腔时我们会用到它）：设 $Y$ 是有限型概形，$i:Y\hookrightarrow X$ 是 $Y$ 到上述形式光滑 $X$ 的闭嵌入。则 $X$ 上支撑在 $Y$ 上（集合论意义下）的 D-模范畴只依赖于 $Y$，而不依赖于 $i$ 与 $X$。事实上，该范畴与 $X$ 上 D-晶体的范畴典范等同。

\bdp{7.11.7} 我们来解释归纳概形情形下什么是微分算子。设 $X$ 是 ind-finite type 的归纳概形。对 $X$ 上 $\mathcal{O}^!$-模 $M$，置
\begin{equation}\tag{348}
\operatorname{Der}(\mathcal{O}_X,M):=\varinjlim\operatorname{Der}(\mathcal{O}_Y,M_{(Y)})=\varinjlim\operatorname{Hom}(\Omega_Y,M_{(Y)}).
\end{equation}
这里 $Y$ 是 $X$ 的闭子概形。把 $\operatorname{Der}(\mathcal{O}_X,M)$ 视为 $X$ 上 $\mathcal{O}^!$-模。类似地，置
\begin{equation}\tag{349}
\mathcal{D}(M)=\operatorname{Diff}(\mathcal{O}_X,M):=\varinjlim\operatorname{Diff}(\mathcal{O}_Y,M_{(Y)}).
\end{equation}
把微分算子层 $\operatorname{Diff}(\mathcal{O}_Y,M_{(Y)})$ 视为 [BB93] 意义下的「微分 $\mathcal{O}_Y$-双模」，即 $Y\times Y$ 上支撑在对角线上的 $\mathcal{O}$-模（集合论意义下）。
% [Tutor Note: 原书 p.306 此处引号两侧都印成右引号（”differential O_Y-bimodule”），底本如此，译文改用中文引号。]

于是 $\mathcal{D}(M)$ 是 $X\times X$ 上支撑在对角线上的 $\mathcal{O}^!$-模（集合论意义下）。按两个 $\mathcal{O}_X$-模结构中的任一个，都可把它视为 $X$ 上 $\mathcal{O}^!$-模。注意 $\mathcal{D}(M)$ 带典范递增滤链 $\mathcal{D}_\cdot(M)$，其中 $\mathcal{D}_i(M)$ 是支撑在对角线第 $i$ 个无穷小邻域上的截面子模；等价地，$\mathcal{D}_i(M)=\varinjlim\operatorname{Diff}_i(\mathcal{O}_Y,M_{(Y)})$ 是阶 $\le i$ 的微分算子子模。有 $\mathcal{D}_0(M)=M$、$\bigcup\mathcal{D}_i(M)=\mathcal{D}(M)$，且 $\operatorname{gr}_\cdot\mathcal{D}(M)$ 上两个 $\mathcal{O}^!$-模结构重合。有显然的嵌入 $\operatorname{Der}(\mathcal{O}_X,M)\subset\mathcal{D}_1(M)$。

\SourcePage{307}

现设 $X$ 形式光滑。下一个命题中，我们按左 $\mathcal{O}$-模结构把 $\mathcal{D}(M)$ 视为 $X$ 上 $\mathcal{O}^!$-模。

\bdp{7.11.8} \emph{命题.} (i) 函子 $\operatorname{Der}(\mathcal{O}_X,\cdot)$、$\mathcal{D}$、$\mathcal{D}_i$ 正合，且与正向极限交换。故存在平坦 $\mathcal{O}^p$-模 $\Theta_X$、$\mathcal{D}_X$，以及 $\mathcal{D}_X$ 的一列平坦子模 $\mathcal{D}_{iX}$ 构成的滤链，使
\[
\operatorname{Der}(\mathcal{O}_X,M)=M\otimes\Theta_X,\qquad\mathcal{D}(M)=M\otimes\mathcal{D}_X,\qquad\mathcal{D}_i(M)=M\otimes\mathcal{D}_{iX}.
\]
(ii) 有典范等同 $\operatorname{gr}_\cdot\mathcal{D}_X=\operatorname{Sym}^\cdot\Theta_X$。

\emph{注.} 7.12.12 中将证明 $\Theta_X$、$\mathcal{D}_X$、$\mathcal{D}_{iX}$ 这些 $\mathcal{O}^p$-模是 Raynaud--Gruson 意义下的 Mittag-Leffler 模（见 7.12.1、7.12.2、7.12.9）。若 $X$ 是 $\aleph_0$-归纳概形，则这些 $\mathcal{O}^p$-模限制到 $X$ 的子概形上局部自由（更精确的陈述见 7.12.13）。

\emph{证明.} (i) 这些函子显然左正合，且与正向极限交换。$\operatorname{Der}(\mathcal{O}_X,\cdot)$ 的右正合性由 $X$ 的形式光滑性推出（把导子 $\mathcal{O}_X\to M$ 标准地解读为态射 $\operatorname{Spec}(\operatorname{Sym}^\cdot M/\operatorname{Sym}^{\ge 2}M)\to X$）。于是得到 $\Theta_X\in M^{pfl}(X,\mathcal{O})$。

(ii) 我们定义典范同构\footnote{在一般情形（基域特征可以非零）下，须把 $\operatorname{Sym}^\cdot$ 换成 $\Gamma^\cdot$：对任意平坦 $A$-模 $P$，把 $\Gamma^n(P)$ 定义为 $P^{\otimes n}$ 中的 $S_n$-不变量。注意（因 $P$ 是投射模的归纳极限）$\Gamma^n(P)$ 平坦，且对任意 $A$-模 $M$ 有 $(M\otimes P^{\otimes n})^{S_n}=M\otimes\Gamma^n(P)$。}
\begin{equation}\tag{350}
\sigma_\cdot:\operatorname{gr}_\cdot\mathcal{D}(M)\xrightarrow{\sim}M\otimes\operatorname{Sym}^\cdot\Theta_X.
\end{equation}
这显然推出命题。

注意对任意 $n\ge 0$，显然的态射 $M\otimes\Theta_X^{\otimes n}\to\varinjlim\operatorname{Hom}(\Omega_Y^{\otimes n},M_{(Y)})$ 是同构（用到 $\Omega_Y$ 凝聚这一事实）。因此 (350) 等价于等同
\begin{equation}\tag{351}
\sigma_n:\operatorname{gr}_n\mathcal{D}(M)\xrightarrow{\sim}\varinjlim\operatorname{Hom}(\operatorname{Sym}^n\Omega_Y,M_{(Y)}).
\end{equation}

\SourcePage{308}

我们的 $\sigma_n$ 是下列映射的归纳极限：
\[
\sigma_{nY}:\operatorname{gr}_n\operatorname{Diff}(\mathcal{O}_Y,M_{(Y)})\to\operatorname{Hom}(\operatorname{Sym}^n\Omega_Y,M_{(Y)}),
\]
它们定义如下。有 $\operatorname{Diff}_n(\mathcal{O}_Y,M_{(Y)})=\operatorname{Hom}_{\mathcal{O}_Y}(\mathcal{O}_{Y\times Y}/\mathcal{I}^{n+1},M_{(Y)})$，其中 $\mathcal{I}\subset\mathcal{O}_{Y\times Y}$ 是对角线理想（源按某个投影映射视为 $\mathcal{O}_Y$-模）。又 $\mathcal{I}/\mathcal{I}^2=\Omega_Y$，故 $\mathcal{I}^n/\mathcal{I}^{n+1}$ 是 $\operatorname{Sym}^n\Omega_Y$ 的商，而 $\sigma_{nY}$ 来自映射 $\operatorname{Sym}^n\Omega_Y\to\mathcal{I}^n/\mathcal{I}^{n+1}\subset\mathcal{O}_{Y\times Y}/\mathcal{I}^{n+1}$。

余下要证 $\sigma_n$ 是同构；可设 $n\ge 1$。显然 $\sigma_{nY}$ 单射，故 $\sigma_n$ 单射。为看出 $\sigma_n$ 满，考虑概形 $Z:=\operatorname{Spec}(\operatorname{Sym}^\cdot\Omega_Y/\operatorname{Sym}^{\ge n+1}\Omega_Y)$。其子概形 $\operatorname{Spec}(\operatorname{Sym}^\cdot\Omega_Y/\operatorname{Sym}^{\ge 2}\Omega_Y)=\operatorname{Spec}(\mathcal{O}_{Y\times Y}/\mathcal{I}^2)\subset Y\times Y\subset Y\times X$ 的嵌入，借 $X$ 的形式光滑性延拓为 $Y$ 上的态射 $i:Z\to Y\times X$。容易看出 $i$ 是闭嵌入。存在闭子概形 $Y'\subset X$，使 $Y\subset Y'$ 且 $Z\subset Y\times Y'$。于是 $Z$ 是 $Y'\times Y'$ 中对角线第 $n$ 个无穷小邻域的子概形。因此得到嵌入 $\operatorname{Hom}(\operatorname{Sym}^n\Omega_Y,M_{(Y)})\subset\operatorname{Hom}_{\mathcal{O}_Y}(\mathcal{O}_Z,M_{(Y)})\subset\operatorname{Diff}_n(\mathcal{O}_{Y'},M_{(Y')})$。把它们与 $\sigma_{nY'}$ 复合，得到嵌入 $\operatorname{Hom}(\operatorname{Sym}^n\Omega_Y,M_{(Y)})\subset\operatorname{Hom}(\operatorname{Sym}^n\Omega_{Y'},M_{(Y')})$。这意味着 $\sigma_n$ 满。\hfill$\Box$

\bdp{7.11.9} 为解释归纳概形上的 D-模是什么，用微分双模的语言比较方便。

设 $X$ 是任意合理归纳概形。$X$ 上的 Diff-双模 $D$（参看 [BB93]）是这样的规则：对任意合理子概形 $Y\subset X$，指派 $Y\times X$ 上一个 $\mathcal{O}^!$-模 $D_Y$，它在 $Y\times X$ 中支撑在（集合论意义下的）对角线 $Y$ 上；对 $Y\subset Y'$ 有等同 $D_{Y'}\otimes\mathcal{O}_Y\xrightarrow{\sim}D_Y$，这些等同按显然意义传递。

Diff-双模的范畴 $M^{di}(X,\mathcal{O})$ 是幺半 $\mathbb{C}$-范畴。即对 $D,D'\in M^{di}(X,\mathcal{O})$，其张量积 $D\otimes D'$ 由
\[
(D\otimes D')_Y:=\varinjlim(D_Y)_{(Y\times Y')}\otimes_{\mathcal{O}_{Y'}}D'_{Y'}
\]
定义。我们的 $\mathcal{O}_X$ 是 $M^{di}(X,\mathcal{O})$ 的单位对象\SourcePage{309}（见下面的注 (i)）。范畴 $M(X,\mathcal{O})$ 是右 $M^{di}(X,\mathcal{O})$-模：对 $\mathcal{O}^!$-模 $M$ 有 $M\otimes D=\varinjlim M_{(Y)}\otimes D_Y$，其中把 $M_{(Y)}\otimes D_Y$ 按 $D_Y$ 上的右 $\mathcal{O}^!$-模结构视为 $X$ 上 $\mathcal{O}^!$-模。
% [Tutor Note: 原书 p.309 作 “a right M^{di}(X,O)-Module”，Module 无故大写；底本排版异常，译文按普通名词处理。]

\emph{注.} (i) $X$ 上 $\mathcal{O}^p$-模等同于支撑在（概形论意义下的）对角线上的微分 $\mathcal{O}_X$-双模。故有幺半范畴的全忠实嵌入 $M^p(X,\mathcal{O})\subset M^{di}(X,\mathcal{O})$。它与 7.11.4、7.11.9 中在 $M(X,\mathcal{O})$ 上的作用相容。
% [Tutor Note: 原书 p.309 此处作 “the Actions on M(X,O)”，Actions 无故大写；底本排版异常，译文按普通名词处理。]

(ii) 遗忘\footnote{忘掉右 $\mathcal{O}$-模结构。}函子 $M^{di}(X,\mathcal{O})\to M^p(X,\mathcal{O})$ 忠实，故可把 Diff-双模看作带额外结构的 $\mathcal{O}^p$-模。称 Diff-双模平坦，指它作为 $\mathcal{O}^p$-模平坦。平坦 Diff-双模的范畴是正合范畴（参看 7.11.3）。

$X$ 上的 Diff-代数指幺半范畴 $M^{di}(X,\mathcal{O})$ 中的幺结合代数 $D$。$X$ 上的 $D^!$-模指 $M(X,\mathcal{O})$ 中的（必为右）$D$-模 $M$。常把这样的 $M$ 简称为 D-模。$D$-模范畴记作 $M(X,D)$，它是阿贝尔范畴。

\emph{注.} (i) 遗忘函子 $M(X,D)\to M(X,\mathcal{O})$ 有左伴随函子，即 $M\mapsto M\otimes D$。

(ii) 范畴 $M^p(X,\mathcal{O})$ 以显然方式成为左 $M^{di}(X,\mathcal{O})$-模。故可考虑 $D^p$-模，即在 $M^p(X,\mathcal{O})$ 中的左 $D$-模。

对 $D\in M^{di}(X,\mathcal{O})$，置 $\Gamma(X,D):=\varprojlim\Gamma(Y\times X,D_Y)$；这是拓扑向量空间。有显然的连续映射 $\Gamma(X,D)\otimes\Gamma(X,D')\to\Gamma(X,D\otimes D')$。对 $M\in M(X,\mathcal{O})$ 有类似的映射 $\Gamma(X,M)\otimes\Gamma(X,D)\to\Gamma(X,M\otimes D)$。因此对 Diff-代数 $D$，$\Gamma(X,D)$ 是拓扑环；且对任意 $D$-模 $M$，向量空间 $\Gamma(X,M)$ 是离散 $\Gamma(X,D)$-模。

设群归纳概形（或任意群「空间」）$K$ 作用在 $X$ 上。$K$ 在 Diff-代数 $D$ 上的弱\footnote{关于强作用见 [BB93]。}作用定义如下。\SourcePage{310}对任意交换代数 $A$，群 $K(A)$ 作用在 $\operatorname{Spec}A\times X$ 上。于是 $K$ 在 $D$ 上的弱作用是下述规则：对每个 $A$，把该作用提升为 $\operatorname{Spec}A\times X$ 上 Diff-代数 $\mathcal{O}_{\operatorname{Spec}A}\boxtimes D$ 上的作用。对任意态射 $A\to A'$，相应作用须按显然方式相容。若 $M$ 是 $D$-模，则 $K$ 在 $M$ 上的弱作用指 $K$ 在 $M$ 上作为 $\mathcal{O}^!$-模的作用（见 7.11.4），使 $D$-作用态射 $M\otimes D\to M$ 与 $K$-作用相容。弱 $K$-等变 $D$-模的范畴记作 $M(K\backslash X,D)$。

\bdp{7.11.10} 对合理 $\aleph_0$-归纳概形 $X$ 上的微分双模与代数，还有更具体的「层论」看法。\footnote{$\aleph_0$ 假设使我们能用拓扑代数代替投射代数；见 7.11.2(i)。}我们分两步解释。
% [Tutor Note: 原书 p.310 的引号两侧都印成右引号（”sheaf-theoretic”），底本如此，译文改用中文引号。]

(i) 设 $X_{red}$ 是概形，即 $X$ 是形式概形\footnote{形式光滑仿射 $\aleph_0$-形式概形（ind-finite type）的描述见 7.12.22、7.12.23。}。于是 $X$ 的底层拓扑空间有定义，$\mathcal{O}_X$ 是拓扑代数的层。任意 Diff-双模 $D$ 给出一个拓扑 $\mathcal{O}_X$-双模的层 $\varprojlim D_{X_\alpha}$，仍滥用记号记作 $D$。它满足下列性质：
\begin{itemize}[label=--]
\item $D$ 上拓扑的基由 $\mathcal{I}\cdot D$ 的闭包构成，其中 $\mathcal{I}\subset\mathcal{O}_X$ 是开理想；该拓扑完备且分离。
\item 商 $D/\mathcal{I}\cdot D$ 是 $X\times X$ 上支撑在（集合论意义下的）对角线处的 $\mathcal{O}^!$-模。
\end{itemize}
显然 $M^{di}(X,\mathcal{O})$ 等价于这类拓扑 $\mathcal{O}_X$-双模层的范畴。注意 $D\otimes D'=D\hat{\otimes}_{\mathcal{O}_X}D'$。因此 $X$ 上的 Diff-代数等同于 $X$ 上拓扑代数的层 $D$，连同代数层的连续态射 $\epsilon:\mathcal{O}_X\to D$，且要求 $D$ 上的 $\mathcal{O}_X$-双模结构满足上述条件。

\SourcePage{311}

$X$ 上的 D-模等同于离散右 D-模的层，它作为 $\mathcal{O}_X$-模拟凝聚（即它是 $X$ 上 $\mathcal{O}^!$-模）。

(ii) 设 $X$ 是任意合理 $\aleph_0$-归纳概形。对合理子概形 $Y\subset X$，记 $Y^\wedge$ 为 $X$ 沿 $Y$ 的完备化。这是 (i) 中的形式概形。对 $X$ 上的 Diff-双模 $D$，令 $D_{Y^\wedge}$ 为 $D$ 到 $Y^\wedge$ 的（$\mathcal{O}^p$-模）拉回。它是 $Y^\wedge$ 上的 Diff-双模，故如 (i) 可看作 $\mathcal{O}_{Y^\wedge}$-双模的层。若 $Y'\subset X$ 是含 $Y$ 的另一个合理子概形，则有 $\mathcal{O}_{Y'^\wedge}$-双模层的连续态射 $D_{Y'^\wedge}\to D_{Y^\wedge}$；它把 $D_{Y^\wedge}$ 与 $D_{Y'^\wedge}$ 关于由 $\mathcal{I}\cdot D_{Y'^\wedge}$ 的闭包生成的拓扑的完备化等同，其中 $\mathcal{I}\subset\mathcal{O}_{Y'^\wedge}$ 是使 $\operatorname{Spec}(\mathcal{O}/\mathcal{I})_{red}=Y_{red}$ 的开理想。这些态射满足显然的传递性。显然 $X$ 上的 Diff-双模等同于这类数据。

因此 $X$ 上的 Diff-代数 $D$ 可看作如下数据：
\begin{itemize}[label=--]
\item 一簇拓扑代数的层 $D_{Y^\wedge}$，同时对任意合理子概形 $Y\subset X$ 定义满足 (i) 中条件的态射 $\epsilon_{Y^\wedge}:\mathcal{O}_{Y^\wedge}\to D_{Y^\wedge}$；
\item 对 $Y\subset Y'$ 有连续态射 $r_{YY'}:D_{Y'^\wedge}\to D_{Y^\wedge}$，它如上把 $D_{Y^\wedge}$ 与 $D_{Y'^\wedge}$ 的完备化等同。要求相容性 $r_{YY'}\epsilon_{Y'^\wedge}=\epsilon_{Y^\wedge}$、$r_{YY''}=r_{YY'}r_{Y'Y''}$。
\end{itemize}
用这种语言描述 D-模，留给读者。

\emph{注.} 对 Diff-代数 $D$，拓扑代数 $\Gamma(X,D)$ 是拓扑代数 $\Gamma(Y,D_{Y^\wedge})$ 的射影极限。

\bdp{7.11.11} \emph{关键例子.} 设 $X$ 是 ind-finite type 的形式光滑归纳概形。考虑 7.11.8(i) 中定义的 $\mathcal{O}^p$-模 $\mathcal{D}_X$。于是对子概形 $Y\subset X$，$\mathcal{O}_Y$-模 $(\mathcal{D}_X)_Y$ 是 $\mathcal{D}(\mathcal{O}_Y):=\varinjlim\operatorname{Diff}(\mathcal{O}_{Y'},\mathcal{O}_Y)$，带其左 $\mathcal{O}_Y$-模结构。我们的 $\mathcal{D}_X$ 带显然的 Diff-双模结构。微分算子的复合使 $\mathcal{D}_X$ 成为 $X$ 上的 Diff-代数。由 7.11.8，$\mathcal{D}_X$ 带典范环滤链 $\mathcal{D}_{iX}$，使 $\operatorname{gr}_\cdot\mathcal{D}_X=\operatorname{Sym}^\cdot\Theta_X$。\SourcePage{312}拓扑代数 $\Gamma(X,\mathcal{D}_X)$ 称为 $X$ 上整体微分算子的环。

$\mathcal{D}_X$-模范畴记作 $M(X,\mathcal{D})$，或简记 $M(X)$。若群「空间」$K$ 作用在 $X$ 上，则 $\mathcal{D}_X$ 带典范弱 $K$-作用（由结构的平移定义）。于是有弱 $K$-等变 D-模的范畴 $\mathcal{M}(K\backslash X,\mathcal{D}_X)=\mathcal{M}(K\backslash X)$。

\emph{扭版本.} 本书正文中还需要扭微分算子环（alias tdo）、这类环的族及其上的模。相应定义是有限维情形下通常定义的直接修改（例如参看 [BB93]）。下面具体描述我们需要的 tdo 特例。

设 $X$ 如上，$\mathcal{L}$ 是 $X$ 上的线丛（见 7.11.3）。

a. 作用在 $\mathcal{L}$ 上的微分算子的 Diff-代数 $\mathcal{D}_{\mathcal{L}}$ 的定义与 $\mathcal{D}_X$ 完全相同，只须在 (349) 中把 $\mathcal{D}(M)$ 换成
\[
\mathcal{D}_{\mathcal{L}}(M)=\operatorname{Diff}(\mathcal{L},M\otimes\mathcal{L}):=\varinjlim\operatorname{Diff}(\mathcal{L}_Y,M_{(Y)}\otimes\mathcal{L}_Y);
\]
命题 7.11.8（及其证明）无须任何改动依然成立。等价地，$\mathcal{D}_{\mathcal{L}}=\mathcal{L}\otimes\mathcal{D}_X\otimes\mathcal{L}^{\otimes-1}$。

b. 如下定义 $X$ 上的 Diff-代数 $\mathcal{D}_{\mathcal{L}^h}$。设 $\pi:X^\sim\to X$ 是对应于 $\mathcal{L}$ 的 $\mathbb{G}_m$-torsor（故 $X^\sim$ 是 $\mathcal{L}$ 去掉零截面）。考虑 $X$ 上的 Diff-代数 $\mathcal{D}^\sim:=\pi_*\mathcal{D}_{X^\sim}$（故对子概形 $Y\subset X$ 有 $(\mathcal{D}^\sim)_Y:=\pi_*((\mathcal{D}_{X^\sim})_{\pi^{-1}Y})$）。$\mathcal{D}_{X^\sim}$ 上的弱 $\mathbb{G}_m$-作用给出 $\mathcal{D}^\sim$ 上的弱 $\mathbb{G}_m$-作用（相对 $X$ 上的平凡 $\mathbb{G}_m$-作用）。我们的 $\mathcal{D}_{\mathcal{L}^h}$ 是 $\mathcal{D}^\sim$ 中 $\mathbb{G}_m$-不变量的子代数。

记 $h$ 为 $\mathcal{D}_{\mathcal{L}^h}$ 中对应于 $-t\frac{d}{dt}\in\operatorname{Lie}\mathbb{G}_m$ 的作用的整体截面。则 $\mathcal{D}_{\mathcal{L}^h}$ 是 $h$ 在 $\mathcal{D}^\sim$ 中的中心化子。注意对任意子概形 $Y\subset X$，$\mathcal{L}_{Y^\wedge}$ 的平凡化（它在 $Y$ 上局部存在）给出等同
\[
\mathcal{D}_{\mathcal{L}^hY^\wedge}\xrightarrow{\sim}\mathcal{D}_{Y^\wedge}\otimes_{\mathbb{C}}\mathbb{C}[h].
\]

\emph{注.} (i) 考虑 $\mathcal{O}^p$-模 $\pi_*(\mathcal{O}_{X^\sim})=\oplus\mathcal{L}^{\otimes n}$。它带 $\mathcal{D}_{\mathcal{L}^h}$ 的作用，且该作用保持分次。$\mathcal{D}_{\mathcal{L}^h}$ 在 $\mathcal{L}^{\otimes n}$ 上的作用把 $\mathcal{D}_{\mathcal{L}^h}/(h-n)\mathcal{D}_{\mathcal{L}^h}$ 与 $\mathcal{D}_{\mathcal{L}^{\otimes n}}$ 等同。

\SourcePage{313}

(ii) 设 $M^\sim$ 是 $X^\sim$ 上弱 $\mathbb{G}_m$-等变的 D-模。置 $M:=(\pi_*M^\sim)^{\mathbb{G}_m}$；它是 $\mathcal{D}_{\mathcal{L}^h}$-模。函子 $\mathcal{M}(\mathbb{G}_m\backslash X^\sim)\to\mathcal{M}(X,\mathcal{D}_{\mathcal{L}^h})$，$M^\sim\mapsto M$，是范畴等价。

\bdp{7.11.12} 我们来解释 D-Ω 复形在归纳概形情形下的相互作用。先定义 Ω-复形。这里设 $X$ 是任意合理归纳概形。

对任意合理子概形 $Y\subset Y'$，有交换 DG 代数的满态射 $\Omega_{Y'}\to\Omega_Y$。$X$ 上的 $\Omega^!$-复形 $F$（或简称 Ω-复形）是这样的规则：对合理子概形 $Y\subset X$ 指派 DG $\Omega_Y$-模 $F_{[Y]}$，并对 $Y\subset Y'$ 给出 $\Omega_{Y'}$-模的态射 $F_{[Y]}\to F_{[Y']}$，这些态射把 $F_{[Y]}$ 与 $i^!_{\Omega YY'}F_{[Y']}:=\operatorname{Hom}_{\Omega_{Y'}}(\Omega_Y,F_{[Y']})$ 等同，且满足显然的传递条件。假定 $F^i_{[Y]}$ 作为 $\mathcal{O}_Y$-模拟凝聚。如 7.11.4，只须考虑 $X_\alpha$，不必考虑所有合理 $Y$。如 7.2.1 的注，这样的 $F$ 等同于一个 $\mathcal{O}^!$-模复形，其微分是阶 $\le 1$ 的微分算子。记 $C(X,\Omega)$ 为 $\Omega^!$-复形的 DG 范畴。

若 $f:Y\to X$ 是归纳概形的可表示拟紧态射（故 $Y=\varinjlim Y_\alpha$，其中 $Y_\alpha:=f^{-1}(X_\alpha)$），则有拉回函子
\[
f^\cdot_\Omega:C(X,\Omega)\to C(Y,\Omega),\qquad f^\cdot_\Omega(F):=\varinjlim\Omega_{Y_\alpha}\otimes_{f^{-1}\Omega_{X_\alpha}}F_\alpha.
\]
若 $f$ 满且形式光滑，则 $f^\cdot_\Omega$ 满足下降性质。

设群「空间」$K$ 作用在 $X$ 上。$K$ 在 $X$ 上 Ω-复形 $F$ 上的作用定义如下：对任意 $g\in K(A)$，把 $g$ 在 $\operatorname{Spec}A\times X$ 上的作用提升为 $\mathcal{O}_{\operatorname{Spec}A}\otimes F\in C(\operatorname{Spec}A\times X,\Omega)$ 上的作用；且须满足显然的相容性。相应范畴记作 $C(K\backslash X,\Omega)$。

\emph{注.} (i) 设 $K$ 是群归纳概形，于是有李代数 $\operatorname{Lie}K$。7.6.4 中用算子 $i_\xi$ 给出的 $K_\Omega$-作用的定义，立即适用于当前情形。$K_\Omega$-等变 Ω-复形的范畴记作 $C(K\backslash X,\Omega)$。

\SourcePage{314}

(ii) 若我们的 $K$ 是仿射群概形，则 $K_\Omega$-等变 Ω-复形等同于这样的 Ω-复形 $F$：它带 $K\times X$ 上 Ω-复形的同构 $m^\cdot_\Omega F=p^\cdot_X F$，且满足通常的条件（见 7.6.5）。

\bdp{7.11.13} 设 $X$ 是 ind-finite type 的形式光滑归纳概形。记 $C(X,\mathcal{D})$ 为 $X$ 上 D-模复形（简称 D-复形）的 DG 范畴。有 DG 函子
\begin{equation}\tag{352}
\mathcal{D}:C(X,\Omega)\to C(X,\mathcal{D}),
\end{equation}
它把 Ω-复形 $F$ 送到 D-复形 $\mathcal{D}F$，其分量为 $(\mathcal{D}F)^n:=\mathcal{D}(F^n)=F^n\otimes\mathcal{D}_X$（见 7.11.8），微分由公式 $d(a):=d_F\circ a$ 定义（这里 $a\in\mathcal{D}(F^n)=\operatorname{Diff}(\mathcal{O}_X,F^n)$）。该函子有右伴随函子
\begin{equation}\tag{353}
\Omega:C(X,\mathcal{D})\to C(X,\Omega),
\end{equation}
它可以具体描述如下。对子概形 $Y\subset X$，有 D-复形 $DR_Y:=\mathcal{D}(\Omega_Y)$。它也是左 DG $\Omega_Y$-模。于是对 D-复形 $M$ 有
\[
\Omega M=\varinjlim\operatorname{Hom}(DR_Y,M)=\bigcup\operatorname{Hom}(DR_Y,M).
\]
引理 7.2.4 连同其证明仍然成立。如 7.2.5，有上同调函子 $H^\cdot_D:C(X,\Omega)\to\mathcal{M}(X)$，$H^\cdot_D(F)=H^\cdot(\mathcal{D}F)$，以及相应的 D-拟同构概念。$\mathcal{D}$、$\Omega$ 的伴随态射是拟同构与 D-拟同构\footnote{D-模的 de Rham 复形不从下界有界，这一事实并不妨碍上述图景。}。

\bdp{7.11.14} 称一个 $\mathcal{O}^!$-复形或 $\mathcal{O}^!$-模具有\emph{拟紧支撑}，若它在某个闭子概形的补集上为零。同一定义适用于 D-复形与 Ω-复形。我们用下标「c」标记相应的范畴。函子 $\mathcal{D}$ 与 $\Omega$ 保持相应的满 DG 子范畴 $C_c(X,\Omega)\subset C(X,\Omega)$、$C_c(X,\mathcal{D})\subset C(X,\mathcal{D})$。

\SourcePage{315}

为确保我们的导出范畴正是所要的（即它们具有好的函子性质），我们额外假设\emph{对角线态射 $X\to X\times X$ 是仿射的}（参看 7.3.1）。例如，只要假设 $X$ 分离就够。

记 $D(X,\mathcal{O})$ 为 $C_c(X,\mathcal{O})$ 的同伦范畴关于拟同构的局部化；这是以 $M_c(X,\mathcal{O})$ 为核心的 t-范畴。类似地定义 $D(X,\mathcal{D})$（假设 $X$ 是 ind-finite type 的形式光滑概形）；这是以 $M_c(X)$ 为核心的 t-范畴。设 $D(X,\Omega)$ 是 $C_c(X,\Omega)$ 的同伦范畴关于 D-拟同构的局部化。函子 $\mathcal{D}$ 与 $\Omega$ 给出 $D(X,\mathcal{D})$ 与 $D(X,\Omega)$ 的典范等同。故照例把这些等同后的范畴简记为 $D(X)$\footnote{为得到以 $\mathcal{M}(X)$ 为核心的 t-范畴，可以考虑那些为具有拟紧支撑的子复形之并的复形；但要保证该范畴有好的函子性质，须假设 $X$ 满足某些额外条件（例如存在形式光滑满态射 $Y\to X$，其中 $Y$ 归纳仿射）。由所有复形构成的范畴具有不好的同调与函子性质。注意通常的补救办法——只考虑从下界有界的 Ω-复形——在这里不起作用（D-模的 de Rham 复形不满足该条件）。}。

称具有拟紧支撑的 $\mathcal{O}^!$-模 $F$ 是\emph{松弛}的，若对任意使得 $F$ 支撑在 $Y^\wedge$ 上的闭子概形 $Y\subset X$，以及 $Y^\wedge$ 上任意平坦 $\mathcal{O}^p$-模 $P$，都有 $H^a(X,P\otimes F)=0$（$a>0$）。称 $\mathcal{O}^!$-、D- 或 Ω-复形 $F$ 松弛，若每个 $\mathcal{O}^!$-模 $F^i$ 都松弛。有下述与 7.3.8 平行的引理：

\bdp{7.11.15} \emph{引理.} i) 对任意 $F'\in C_c(X,\Omega)$，存在 D-拟同构 $F'\to F$，使 $F$ 松弛，且 $F$ 与 $F'$ 的支撑相同。
% [Tutor Note: 原书 p.315 此项标号印作 i)，下一项印作 (ii)，底本如此，照录。]

(ii) 若 $f:X\to X'$ 是归纳概形的形式光滑仿射态射，则函子
\[
f^\cdot_\Omega:C_c(X',\Omega)\to C_c(X,\Omega),\qquad f^\cdot:C_c(X,\Omega)\to C_c(X',\Omega)
\]
把松弛复形送到松弛复形。

% 尾部接 p.316（承接 7.11.15 的 (iii)）
